%% The following is a directive for TeXShop to indicate the main file
%%!TEX root = diss.tex

%================================
\chapter{Generalizations}
\label{ch:General}
%================================

%================================
\section{Introduction}
%================================
We will push the operator notation further to see what kinds of lattice walk models we can handle. Some of the models have more general step sets while others are set in more complicated regions. For example, Definition 1.1 of \cite{mishna2009classifying} has a functional equation for the non-interacting lattice walks on the quarter plane, which is valid for models with $\mc S$ containing negative coordinates of size at most $\um 1$. Similarly, \cite{kauers2014walks} covers walks in the quarter plane with multiple copies of small steps. We start with the following definition for a lattice walk model in $\ze^d$. Recall the definition from \autoref{ch:Introduction}.
\begin{definition}(MOVE) Let $\mc{D}_0$ and $\mc{D}$ be regions in $\ze^d$, with $\mc{D}_0 \subseteq \mc D$, and let $\mc S$ be a set of $d$-tuples in $\ze^d$. A triple $\mc{M} = (\mc{S}, \mc{D},\mc{D}_0)$ is a lattice walk model with steps $\mc{S}$, living in the region $\mc{D}$ and interacting with region $\mc{D}_0$. We may impose a starting location within $\mc D$ later on.
\end{definition}
In other words, given a starting location in $\mc D$ we can evolve lattice walks in the region $\mc D$ with steps from $\mc{S}$. We also have some special locations $\mc{D}_0$ within $\mc{D}$, and we wish to count the interactions of the lattice walks with points in this region. The interactions may be further specified later on.

\begin{figure}[H]
    \centering
    \includegraphics[scale = 0.5]{placeholder}
    \caption{General regions in $\ze^2$.}
\end{figure}

For example, we may consider a semi-infinite strip on the plane as in \cite{brak2005directed}, \cite{wong2014confining} and \cite{owczarek2012enumeration}. In \cite{niederhausen2005note} the author considers lattice walks in an octant with interactions with the diagonal.

\begin{definition} Let $\mc{M}$ be a lattice walk model and $Q$ its associated generating function counting lattice walks' lengths with $t$, ending locations with $\row xd$, and interactions with $\mc{D}_0$ with some multi-index variable $a$. We write
\begin{align}
	Q = Q(\row xd,a;t) = \sum_{k=0}^{\infty} q_k(\row xd, a)t^k\,,
\end{align}
where $q_k$ are Laurent poynomials in the $x$ and $a$ variables. 
\end{definition}
\begin{definition} Suppose that $\Gamma$ is an operator involving the $x$ and $a$ variables but not $t$. We shall apply $\Gamma$ on power series $F = \sum_{k=0}^\infty f_k t^k$ in the following way:
\begin{align}
	\Gamma F = \sum_{k=0}^\infty \Gamma (f_k) t^k\,.
\end{align}
This is to avoid problems in expanding $F$ in Laurent series in some of the $x$ variables.
\end{definition}

\begin{definition} Let $\mc{M}$ be a given model. We denote with $S$ the Laurent polynomial (or series) keeping inventory of the steps $\mc{S}$. Let $K = 1-tS$ be the Kernel of $\mc{M}$. Let $E$ be an operator that eliminates monomials $x_{[d]}^P$ when $P$ is not in $\mc{D}$. Let $L$ be an operator that marks (or labels) monomials $x_{[d]}^P$ for $P \in \mc{D}_0$ in a specified way in the model. Note that $L$ is invertible with inverse $L^{\um 1}$ that is given by marking with $1/a$ instead.
\end{definition}

This operator notation allows us to write down functional equations for $Q$ really quickly.
\begin{theorem} Let $\mc{M}$ be a lattice walk model and $Q$ its associated generating function with an initial term $q_0 = q$. Then $Q$ satisfies the following functional equations
\begin{align}
	(1-tLES)Q &= q\\
	(L^{\um 1}-tES)Q &= L^{\um 1}(q) = \alpha
\end{align}
\end{theorem}
\begin{proof} The second functional equation is obtained by applying $L^{\um 1}$ on the first one, so it only remains to show that the first functional equation holds. We do this via iteration. We start with $Q$ and multiply it with $tS$. This gives $tSQ$ or an enumeration of all walks extended by a single step. Here the power of $t$ is raised by one to keep track of this extension. Then we apply $E$ to this result to eliminate all lattice walks that venture outside of the region $\mc{D}$. Finally, we mark those walks that have arrived at locations in $\mc{D}_0$ by applying $L$ to $E(tSQ)$. The resulting expression is almost $Q$, but one term is missing - namely, $q_0 = q$. That is, we have transformed $Q$ to itself in a nontrivial way and we have the equation
\begin{align}
	Q = q + LE(tSQ) = q + tLE(SQ)\,,
\end{align}
and to complete the proof, we rewrite the equation to the form at the start of the theorem.
\end{proof}

\begin{example} Also, we can consider a different approach where, instead of removing walks that venture outside of $\mc{D}$, we just mark them with some variable $z$ for later deletion. Suppose that we are given some model $\mc{M}$ and we want to frame the analysis in this new way. We have a new lattice walk model $\wcal{M}$ with $\wcal{D} = \ze^d$ and $\wcal{D}_0 = \mc{D}_0 \cup \mc{D}^c$, where we mark monomials in $\mc{D}^c$ with $z$ and those in $\mc{D}_0$ in the same way as in model $\mc{M}$. Here $\widetilde{E}$ is the identity operator since we do not remove any walks. We have the functional equations
\begin{align}
(1-t\widetilde{L} S)\widetilde Q &= q\,,\\
(\widetilde{L}^{\um 1}-tS) Q &= \widetilde{\alpha}\,.
\end{align}
It is not clear how one would proceed and if this would be a better way to handle the problem. The upside is that the operators are invertible, but the downside is that we now have more extracted terms in the functional equation, since the lattice walks can now go anywhere.
\end{example}

Next, we can talk about some of the abstract properties of the operators $L$ and $E$. The operator $E$ has the eigenvalues 1 or 0 for monomials $x_{[d]}^P$ where $P \in \mc{D}$ or $P \in \mc{D}^c$ respectively. The latter characterizes the monomials contained in $\ker E$. Similarly, $L$ has eigenvalues 1 for when $P \in \mc{D}_0^c$, and the marking monomial $a$ for when $P\in \mc{D}_0$. Here $a$ varies from point to point. Also, $L^{\um 1}$ has the reciprocal eigenvalues of $L$.

The task for the rest of the chapter is to come up with representations for the operators involved. In general, this should not be an easy task because if we could express arbitrary operators $L$ and $E$ with a small number of coefficient extractions, we would also be able to use this scheme to compress data without loss, and this is impossible in general.

%================================
\section{Reformulations of Problems}
%================================

\subsection{Using marking operators instead of extraction operators}
An initially promising idea is to use a marking operator instead of an extraction operator with the goal of removing those marked walks at a later stage. The added benefit is that we would be working with invertible operators throughout the majority of the analysis. The downside is that now the walks will not be constrained and $Q$ will no longer be a power series in the $x$ variables. We may get further problems. Instead of $E = 1- \ext{\ol x}$ we take $\tilde E = 1 + (a_{\um 1}-1)\ext{x}$ with the idea that if we set $a_{\um 1} = 0$ we will obtain $E$. To this end the functional equation is
\begin{align}
	(1-tL\tilde E S)\tilde Q = 1
\end{align}
which after inversion turns into
\begin{align}
	(\tilde E^{\um 1} L^{\um 1} -tS) \tilde Q = \tilde E^{\um 1} L^{\um 1} 1 = \alpha
\end{align}
One benefit of this is that we do not have a mixing between the $E$, $L$ operators and $S$. The downside is we have more $Q$ extractions on their own. We can write this in kernel form as
\begin{align}
	K \tilde Q = \alpha + (1 - L^{\um 1}\tilde E^{\um 1}) \tilde Q
\end{align}
where the operator $1- L^{\um 1}\tilde E^{\um 1}$ is composed entirely of coefficient extractions.
\subsection{Considering a symmetrical super-problem}
An extension of the previous idea is to symmetrize the problem by also marking parts on the other side of the origin. Instead of just marking walks on the axes and on lines where one coordinate is $-1$ we can also mark walks on lines with one coordinate $+1$. This leads to a symmetrical arrangement of marking operations on both sides of the origin. We have the equation
\begin{align}
	(1- tLS)\tilde Q = (1- t L_{\um 1}L_0L_1 S)\tilde Q = 1
\end{align}
We can skip the marking of the axes with $L_0$ if we do not wish to investigate wall interactions. In general, we may mark more locations in space so we can combine all of these operators into one called $L$. Then we have
\begin{align}
	(1- tLS)\tilde Q = 1
\end{align}
or in kernel forms
\begin{align}
	K\tilde Q = L^{\um 1}(1) + (1- L^{\um 1})\tilde Q
\end{align}

%================================
\section{Marking Operators}
%================================
The example in the previous section shows that we can talk about marking operators on their own without considering operators $E$. Alternatively, the operators $E$ are just marking operators that apply a label $a = 0$. Therefore, we can combine the discussions for operators $L$ and $E$. Whenever invertibility is concerned, the discussion will apply to operators $L$ only. Suppose that we wish to mark points $p \in P\subseteq \ze^d$ with labels $a_p$ via an operator $\Gamma_P$. It can be written as
\begin{align}
	\Gamma_P = 1 + \sum_{p\in P}(a_p-1) \ext{x_{[d]}^p}\,.
\end{align}
The form $a_p-1$ follows from the fact that we have the original term and we need to remove it before marking with $a_p$. In the case where we mark points for deletion, $P = {\mc{D}}^c$ and $a_p = a$. In the models we have investigated so far we saw that we do not really need the whole operator but just some terms near the boundaries.
\begin{example} If we consider lattice walks with small steps on the half-line in one dimension, we have the operator
\begin{align}
	E = 1 - \sum_{n=1}^\infty \ext{\ol x_1^n}\,,
\end{align}
but since we only have small steps, there are no walks that go past $\ol x_1$ in the negative direction. Therefore, we can use the operator $E' = 1- \ext{\ol x_1}$ to the same effect instead. Although the operators are different in form, they are essentially the same with regards to the problem at hand. We say they are model-equivalent.
\end{example}
To make things easier, we may think of the lattice walk models as living completely within the sumset $\mc{S}+\mc{D}$ of points that are combinations of points in $\mc{D}$ and steps $\mc{S}$. This is how far walks from the model can reach from within $\mc{D}$ by taking a single step.
\begin{definition} Given sets $A$ and $B$, we let their sumset be
\begin{align}
	A+ B = \{a+b\,|\,a\in A,\, b\in B\}\,.
\end{align}
\end{definition}
\begin{definition} Given a lattice walk model $\mc{M}$ with domain $\mc{D}$ and step set $\mc{S}$, define the reach of model $\mc{M}$ to be the sumset $\mc{D}+\mc{S}$.
\end{definition}
With this in mind, we can define a notion of equivalence for operators.
\begin{definition} Let $\mc{M}$ be a lattice walk model in $\ze^d$ and let $\Gamma$ and $\Gamma'$ be two operators acting on $\ze^d$. Then $\Gamma \equiv_{\mc{M}} \Gamma'$ if and only if $\Gamma$ and $\Gamma'$ act identically on monomials from the reach of $\mc{M}$. That is, we call these operators model-equivalent if and only if $\Gamma  x^P = \Gamma' x^P$ for all $P\in \mc{D}+\mc{S}$.
\end{definition}
In the preceding example $E\equiv E'$ since they act the same way on monomials from the reach of $\mc{M}$, that is, monomials $x_1^n$ for $n\ge \um 1$. Similarly, in the two-dimensional problem we have an equivalence between the operator $E$ and the operator $E' = (1-\ext{\ol x_1})(1-\ext{\ol x_2})$.
\begin{align}
	\Gamma = 1 + (a-1)\sum_{k=1}^{\infty} \ext{\ol x^k} = 1 + (a-1)\sum_{k=1}^{n} \ext{\ol x^k}
\end{align}
When the boundary of $\mc{D}$ is complicated, we may not find nice model-equivalent operators, but there are some cases where we can say something. In problems featuring half-planes, we have an operator $E$ with infinitely many terms, and we find that there is a model-equivalent operator $E'$ that is only comprised of finitely many terms. We will investigate which operators are equivalent to finite ones.
\begin{definition} Given a model $\mc{M}$, a coefficient extraction operator $\Gamma$ in the variables $\row xd$ is finite if it is equivalent to an operator $\Gamma'$ comprised of finitely many coefficient extraction terms. That is,
\begin{align}
	\Gamma \equiv_{\mc M} \Gamma' = \sum_{\substack{I\subseteq [d]\\ J \subseteq \ze^d}}a_{I,J} \ext{x_I^J}
\end{align}
where each $|J| < \infty$. Also, $\Gamma'$ is infinite if it is not of finite.
\end{definition}
For most purposes we will not distinguish between $\Gamma$ and $\Gamma'$. A natural question is to consider which models $\mc{M}$ lead to finite operators. This very much depends on the region $\mc D$ and the nature of the marking of region $\mc{D}_0$. We have some cases where it is easy to say.

\begin{lemma} Let $\mc D$ be a $d$-dimensional box in $\ze^d$ and $|\mc S| < \infty$. Then the operator $\Gamma = LE$ is finite.
\end{lemma}
\begin{proof} The operator $L$ is finite since it only needs to mark finitely many points. We can find $E$ explicitly. Let the box have $x_k$ coordinates $[a_k,b_k]$. Suppose that the steps in $\mc S$ have $x_k$ coordinates of absolute size at most $n_k$. If $n_k = 0$, the walks do not change their $x_k$ coordinate and the walks do not venture past the boundaries in that coordinate. When $n_k > 0$ we need to remove all monomials $x^P$ with $P_k \in [a_k-n_k,a_k-1]\cup [b_k+1,b_k+n_k]$. 

\begin{figure}[H]
    \centering
    \includegraphics[scale = 0.5]{placeholder}
    \caption{Finite region enclosed in box $B$.}
    \label{fig:finitebox}
\end{figure}

We can then write the operator as
\begin{align}
	E_k = 1 - \sum_{j=1}^{n_k} \ext{x_k^{a_k-j}}+\ext{x_k^{b_k+j}}\,,
\end{align}
which is equal to $E_k = 1$ when $n = 0$. Applying $\row Ed$ sequentially ensures that we remove all monomials outside of $\mc D$. That is, $E = \mon Ed$. Since each of the operators is finite, their composition $\Gamma = LE$ is finite as well.
\end{proof}
\begin{theorem} Let $\mc D$ be a finite region in $\ze^d$ and $|\mc S| <\infty$. Then the operator $\Gamma = LE$ is finite.
\end{theorem}
\begin{proof} The region $\mc D$ lies within a bounding box $B$, so we may write $E = E' E''$, where $E' = \mon Ed$ from the lemma and $E''$ is the operator removing all monomials outside of $\mc D$ but inside the interior of $B$. $E''$ is finite since there are only finitely many points in that region. Since we are marking only finitely many points $\mc D_0 \subseteq \mc D$, $L$ is finite, and finally, $\Gamma$ is finite.
\end{proof}

Also, we may handle semi-infinite regions if we are only marking finitely many points $\mc D_0$. The analysis is the same as above except that now one or more of $a_k, b_k$ may be infinite, and in this case we actually lose parts of $\Gamma_{B^c}$ and the operator is still finite.

\begin{definition} A lattice walk model with finitely many steps is finite if its extraction and marking operators $E$ and $L$ are finite, and its starting condition $q$ has a finite degree. Otherwise, we say the model is infinite.
\end{definition}
\begin{example} All the models so far has been finite, but not all models are finite. Consider a one-dimensional model on the whole line where we mark points $x_1^i$ with monomials $a_i$ using operator $L$. Then all equivalent marking operators $L'$ would have to have a term like $(a_i-1)\ext{x_1^i}$; otherwise, it would not act the same way on monomials in the model's reach $\ze$. Therefore, $L$ is not equivalent to a finite operator.
\end{example}
Finite models are nicer because they have functional equations with finitely many terms and there is a chance that they will be amenable to the sort of analysis we have performed for quarter-plane walk models in the past. What follows is a rather easy application of the theory so far.

%=======================================
\section{Kernel Forms}
%=======================================
We can write out the functional equations in a way that highlights the kernel $K$ in the equation. We start out as follows
\begin{align}
	(1-tLES)Q &= 1\\
	(1-t(LE-1+1)S)Q &=1\\
	KQ + t(1-LE)SQ &= 1\,.
\end{align}
This is the form we see in the master functional equation. Another way we could have expanded is as follows
\begin{align}
\begin{split}
	(1-tLES)Q &= 1\\
	(1+LE(1-tS-1))Q &= 1\\
	LEKQ+(1-LE)Q &= 1\\
	LEKQ+(1-L)Q &= 1
\end{split}	
\end{align}
where $EQ = Q$ since $Q$ is already of a valid form. We can go with both ways at the same time as well.
\begin{align}\begin{split}
	KQ + t(1-LE)SQ &= 1\\
	KQ + (1-LE)(1-tS-1)Q &=1\\
	KQ + (1-LE)KQ - (1-L)Q &= 1\,.
\end{split}\end{align}
We can apply the inverse $L^{\um 1}$ to those equations too, in particular on the second one.
\begin{align}
\begin{split}
	L^{\um 1}(LEKQ+(1-LE)Q) &= 1
\end{split}	
\end{align}

The second form of the functional equation involving $L^{\um 1}$ should have less terms than the first one by virtue of separating $L$ and $E$ from one another. Using extraction operators, we can constrain lattice walks to lie within a finite region, semi-infinite strip, a semi-infinite sector or any other discrete shape (with diminishing returns for more complicated shapes, of course). Furthermore, using marking operators, we can easily track the behavior of lattice walks as they cross or interact with boundaries or interact with regions within. For example, we may consider walks interacting with the boundary of a box within some region in space. 

%================================
\section{Examples of Regions}
%================================

ZZPICTUREZZ

ZZPICTUREZZ

\begin{example}[Self-avoiding walks]

We can handle discrete self-avoiding lattice walks with this theory. If we pick a marking operator $\Gamma$ that marks every location on the plane, this will allow us to track the lattice walks' history. All we need to do in the end is remove walks that have stepped at some location more than once. However, we will investigate this problem in a bit more detail as it serves as an example of the theory. To use the generating function iteration properly we need a sort of limiting argument. For the first $n$ steps we "only" require an operator covering the square region of $(2n+1)^2$ lattice points. The overall operator is defined through a limiting process of applying the full operator (or truncated version) on a finite Laurent polynomial. Let
\begin{align}
\Gamma_{i,j} = 1+b_{i,j}\ext{x_1^ix_2^j}
\end{align}
where as usual $b_{i,j} = a_{i,j}-1$. This operator takes a generating function and marks the monomial $x_1^ix_2^j$. We will be using a lot of these together, so we let
\begin{align}
\Gamma_n = \prod_{-n\leq i,j\leq n}\Gamma_{i,j}
\end{align}
To find a functional equation for SAWs on the plane we perform an iteration. Let $Q$ be the generating function for walks with history, that is, they have a mark for each coordinate they pass through, and let $Q^{\leq n}$ be the truncated version of this generating function. It is given by
\begin{align}
Q^{\leq n} = \sum_{k=0}^{n} Q_k t^k
\end{align}
where the initial condition is $Q_0 = a_{0,0}$. Starting from $Q_n$ we can find $Q_{n+1}$ by iterating with the step set $S$, marking with $\Gamma_{n+1}$ as follows
\begin{align}
	Q_{n+1} = \Gamma_{n+1}SQ_n
\end{align}
Using a computer we can easily iterate this to add more and more steps. We can also obtain $Q^{\leq n+1}$ from $Q^{\leq n}$ since
\begin{align}\begin{split}
	t\Gamma_{n+1}SQ^{\leq n} & = \sum_{k=0}^n \Gamma_{n+1}SQ_kt^{k+1}= \sum_{k=0}^n \Gamma_{k+1}SQ_kt^{k+1} = \sum_{k=0}^n Q_{k+1}t^{k+1} = Q^{\leq n+1}-Q_0
\end{split}\end{align}
or in other words
\begin{align}
	 Q^{\leq n+1} = a_{0,0} + t\Gamma_{n+1}SQ^{\leq n}
\end{align}
In the interest of finding a functional equation we would like to take the limiit of this as $n \to \infty$. For this we define $\Gamma$ to be the operator marking all monomials $x_1^ix_2^j$, acting in the following way
\begin{align}
	\Gamma x^i y^j = \Gamma_{\max(|i|,|j|)} x^i y^j
\end{align}
and extended by linearity. The limiting process is
\begin{align}
	Q = a_{0,0} + t \Gamma SQ
\end{align}
We finally have the functional equation for lattice walks with full history
\begin{align}
	(1 -t\Gamma S)Q = a_{0,0}.
\end{align}
As in the example of an infinte marking operator, it stands to reason that $\Gamma$ is of infinite type as well. From the equation we can retrieve the generating function $W$ for the SAWs by removing all walks with power of mark $a_{i,j}$ greater than 1. For this we need an operator that removes terms with more than one copy of a given label. Let
\begin{align}
	R_{i,j} &= \ext{a_{i,j}^0} + \ext{a_{i,j}}\\
	R_n &= \prod_{-n\leq i,j \leq n} R_{i,j}\\
	R &= \lim_{n\to \infty} R_n
\end{align}
where the limit is understood as in the example of $\Gamma$ above. The solution for $Q$ can also be restated as
\begin{align}
	Q = \sum_{k=0}^{\infty} t^k (\Gamma S)^{k}(a_{0,0}) = \sum_{k=0}^{\infty} t^k (\Gamma_k S)^{k}(a_{0,0}) 
\end{align}
We can either redo the analysis, applying $R$ at every iteration, or we can apply $R$ to $Q$ and find
\begin{align}
	W = RQ = \sum_{k=0}^{\infty} t^k R(\Gamma S)^{k}(a_{0,0}) = \sum_{k=0}^{\infty} t^k (R\Gamma S)^{k}(a_{0,0}) 
\end{align}
We can write this as a functional equation
\begin{align}
	\mc{L} W = (1-t\Sigma)W = (1-tR\Gamma S)W = a_{0,0}.
\end{align}
It is not clear if the operator $R\Gamma$ is finite. Now it is enough to set all the remaining labels to 1 and we will find the full generating function of SAWs with steps from $\mc{S}$.

The discussion lends itself nicely to generalization where we have $d$ variables and operators $\Gamma_{i_1,\ldots,l_d}$, $R_{i_1,\ldots,i_d}$, and so on, defined analogously. The eigenpair discussion works out in the same way with eigenfunctions for $\lambda = 1$
\begin{align}
	u = \frac{a_{i_1,\ldots,i_d}^2 a_{j_1,\ldots,j_d}^m}{a_{k_1,\ldots,k_d}S}x_1^{r_1}\cdots x_d^{r_d} t^n
\end{align}
Also, we can constrain the SAWs on a part of $\ze^d$ with an extraction operator $E$ and mark special locations with an operator $L$. The functional equation is
\begin{align}
	(1- tLE\Sigma)W = a_{0,\ldots,0}
\end{align}
for arbitrary number of dimensions, with regional constraints and marking of special locations.
\end{example}

%==============================================================================
\section{Further Research Directions}
%==============================================================================

So far we have explored using operators to simplify the derivation of functional equations pertaining to lattice walk models on $\ze^d$, and to simplify the first steps in solutions to these problems. It would be great to know if we can further extend the work into a territory like Sturm-Liouville theory, where we may investigate the spectra of operators and find alternative solutions to the functional equations via eigenfunction expansions. For this to work, we need a suitable inner product that will allow us to answer questions like are the operators self-adjoint. Unfortunately, we have not found such suitable inner products apart from those on symmetric functions. While they seem promising, we have not been able to go far enough with them. In this chapter we will provide the start of such analysis which may be continued in future research.

The operators we have seen so far are marking operators, like $L$ and $E$, in combination of multiplication by another function, like $\Gamma S$, or a combination of operators, like $L^{\um 1} - tES$. In this order, we will find out their properties - eigenvalues, eigenvectors, kernels and invertibility.


%=======================================
\subsection{Marking Operators}
%=======================================

The building block of the operators are coefficient extractions $[x_I^J]$ and term extractions $\ext{x_I^J}$. We will use the following notation for brevity.
\begin{definition} (MOVE) For sets $I$ and $J$ of equal size, we define the term extraction operator
\begin{align}
	[[x_I^J]] \equiv \ext{x_I^J}\,.
\end{align}
\end{definition}
The first question we will answer is what happens when we compose term extraction operators. Suppose that we want to perform the extractions with monomials $x_I^P$ and $x_J^Q$. Then the composition contains three extractions - one for the $x$ powers appearing in $x_I^P$ but not in $x_J^Q$, one for the reversed situation, and one for the common terms between the two monomials. However, the resulting composition will be the zero operator if the power sequences for the common variables do not match.
\begin{example} We compute
\begin{align}
	[[x_{1,2}^{3,4}]][[x_{2,3}^{4,1}]] = [[x_{1,2,3}^{3,4,1}]]
\end{align}
and
\begin{align}
	[[x_{1,2}^{3,2}]][[x_{2,3}^{4,1}]] = 0\,.
\end{align}
\end{example}
The next highly-technical lemma makes the composition precise.
\begin{lemma} (REWRITE) Let $I,\, J,\, P,\, Q$ be sets and $K = I\cap J$, $I' = I\backslash K$, $J' = J \backslash K$. Then let $P' = P|_{I'}$, $Q' = Q|_{J'}$ and $P'' = P\backslash P'$, $Q'' = Q\backslash Q'$. The composition of two term extractions is
\begin{align}
	[[x_I^P]][[x_J^Q]] = \delta(P'',Q'')[[x_{I'}^{P'}x_{J'}^{Q'}x_K^{P''}]]\,.
\end{align}
That is, the resulting operator is zero if the power sequences $P$ and $Q$ do not match over the variables common to the two extractions. Furthermore, the two term extraction operators commute.
\end{lemma}

Combinatorially speaking, if a marking operator marks several different monomials, we may express this operator as a product of marking operators, each marking a single monomial. We can make this notion precise with the algebraic form of the marking operators.
\begin{definition} The operator $\Gamma$ marking a monomial $x_I^J$ with label ``$a$" is given by
\begin{align}
	\Gamma = 1  +(a-1)\ext{x_I^J} = 1 + b[[x_I^J]]\,.
\end{align}
A finite marking operator is one that marks finitely-many monomials. An infinite marking operator is one that is not finite.
\end{definition}
\begin{lemma} Let $\alpha$ and $\beta$ be the labels for two different monomials $x_I^P$ and $x_J^Q$ marked by $\Gamma$. Then
\begin{align}
	\Gamma = 1 + (\alpha-1)[[x_I^P]]+(\beta-1)[[x_J^Q]] = (1+ (\alpha-1)[[x_I^P]])(1+ (\beta-1)[[x_J^Q]])\,.
\end{align}
\end{lemma}
\begin{corollary} Marking operators commute.
\end{corollary}
\begin{proof} This can be seen from their combinatorial nature. Since multiplication is commutative and the marking is done via multiplication, the order of applying the marking operators does not matter. Algebraically, this is a consequence of the commutativity the term extraction operators.
\end{proof}

Combinatorially speaking, we can easily find some of the eigenvalues and eigenfunctions by looking at single monomials. If a monomial is not marked by $\Gamma$, then it is an eigenfunction with eigenvalue 1, since $\Gamma$ does not change it; if it is, then it is an eigenfunction with eigenvalue equal to the label given by $\Gamma$. These are all the polynomial eigenfunctions. This does not mean there cannot be more eigenfunctions among those expressible by power/Laurent series and beyond.
\begin{prop} Let $\Gamma$ be an operator marking a single monomial $x_I^J$ with label $a$. Let $\lambda$ and $u_\lambda$ be an eigenpair. Then $\lambda = a$ with $u_\lambda = x_I^J$, or $\lambda = 1$ with $u_\lambda$ a function not containing $x_I^J$, in the formal power series sense. In particular, every monomial other than $x_I^J$ is an eigenfunction with eigenvalue 1.
\end{prop}
\begin{proof}
Let $u_\lambda = u' + u''x_I^J$ where $u \neq 0$ and $u''$ does not share variables with $x_I^J$. Then
\begin{align}
	\Gamma(u) = u' + au''x_I^J  = \lambda (u' + u''x_I^J)\,,
\end{align}
which we may write as
\begin{align}
	(1-\lambda) u' + (a-\lambda)u''x_I^J  = 0\,.
\end{align}
We have two cases.
\begin{itemize}
\item If $u' = 0$, then $u'' \neq 0$ and $\lambda = a$. The eigenfunction is $x_I^J$.
\item If $u' \neq 0$ and there is no term in $u''x_I^J$ that can compesate for this, then $\lambda = 1$ and $u'' = 0$.
\end{itemize}
\end{proof}

\begin{corollary} The kernel of a finite marking operator $\Gamma$ is generated by the monomials marked with label 0.
\end{corollary}

The proofs can be extended for finite marking operators. It is not entirely clear how we may prove a similar result for operators marking infinitely many monomials.

%=======================================
\subsection{Marking Operators With Multiplication}
%=======================================
We call an operator inhomogeneous if it consists of a homogeneous operator $\Gamma$ and a (possibly infinite) laurent polynomial $S$, and acts in one of the following ways
 \begin{align}
	\Gamma S(u) = \Gamma(Su)\quad\text{or}\quad S\,\Gamma u = S\cdot \Gamma u\,.
 \end{align}
This kind of operator is important because it can be found in the functional equations for lattice walk models. To find their eigenpairs, we first show that the definitions are mostly equivalent.
\begin{prop} The two operators $\Gamma S$ and $S\, \Gamma$ have the same eigenvalues with eigenfunctions differing by a multiple of $S$. That is, $(\lambda,u)$ is an eigenpair for $\Gamma S$ if and only if $(\lambda, Su)$ is an eigenpair for $S\,\Gamma$. In particular, if $\lambda = 0$, we have
\begin{align}
	\ker(\Gamma S) = \ol S \ker(\Gamma)\,.
\end{align}
\end{prop}
\begin{proof} Note that
\begin{align}
	\Gamma Su = \lambda u \quad \Longleftrightarrow \quad S\,\Gamma (Su) = \lambda (Su)\,.
\end{align}
Therefore, the operators have the same eigenvalues, and they have the eigenfunctions $u$ and $Su$ respectively. Now, if we pick $\lambda = 0$ to find the kernels of the operators, we find that $u \in \ker(\Gamma S)$ if and only if $Su \in \ker(S\,\Gamma) = \ker(\Gamma)$, and this completes the proof.
\end{proof}
As the next example shows, it may not be easy to find the eigenpairs for such operators.
\begin{example}
Let $\Gamma$ be a marking operator and let $u$ be an eigenfunction of $\Gamma S$, where 
\begin{align}
	S = x\com u = \sum_k u_k x^k\,.
\end{align}
Computing $(\Gamma S-\lambda I)u$ leads to the system of equations
\begin{align}
	\lambda u_n &= a u_{n-1}\,,\\
	\lambda u_k &= u_{k-1}, \quad \text{if } k \neq n\,.
\end{align}
We have a doubly-infinite recursion, so we can pick an arbitrary initial condition. If $u_n = 0$, the rest of the terms will be zero as well, so we must have $u_n \neq 0$. Set $u_n = 1$. Then, via induction, we can find
\begin{align}
	u_{n+k} &= 1/\lambda^{k}, \quad \text{for $k \ge 0$}\,,\\
	u_{n-k} &= \lambda^{k}/a^k, \quad \text{for $k \ge 1$}\,;
\end{align}
therefore, we have an eigenfunction for each $\lambda \neq 0$. Computing $u(x)$ formally leads to
\begin{align}
	u(x) &=-\frac{\lambda b x^{n+1}}{(\lambda-x)(\lambda-ax)}\,.
\end{align}
\end{example}
The situation becomes more complicated for more complicated functions $S$.
%=======================================
\subsection{Lattice Walks Operator}
%=======================================
In this section we explore the operator $\mc{L} = 1 - t\Gamma S$ for a step Laurent polynomial $S$ and some marking operator $\Gamma$. Let $(u, \lambda)$ be an eigenpair for $\mc{L}$. Then we calculate
\begin{align}\begin{split}
	(\mc{L}-\lambda I) u &= u - \lambda u -t\Gamma Su\\
	t \Gamma S u &= (1-\lambda)u = \mu u\,.
\end{split}\end{align}
If $\lambda = 1$, we would have to compute the kernel of $\Gamma S$, which we did in the previous section. The eigenfunctions of $\mc{L}$ for eigenvalue $\lambda = 1$ are the functions $u \in \ol S\ker{\Gamma}$. Let $u$ be given by a finite order Laurent series
\begin{align}
	u = u_n t^n + u_{n+1}t^{n+1} + \cdots\,,
\end{align}
where $u_n \neq 0$. Then we can rewrite the functional equation as
\begin{align}
	t^{-n}(1-t\Gamma S)u &= \lambda t^{-n}u\\
	(\lambda -1) t^{-n}u &= -t^{1-n}\Gamma Su\,.
\end{align}
Letting $t \to 0$ in this equation yields
\begin{align}
	(\lambda - 1)u_n = -\lim_{t\to 0} t^{1-n}\Gamma S u = -\Gamma S \lim_{t\to 0} (t^{1-n} u) = 0\,.
\end{align}
Since $u_n \neq 0$, this implies $\lambda = 1$, and again we see that $u \in \ol S\ker \Gamma$. Since $S$ is a Laurent polynomial, the resulting functions will indeed be finite order Laurent series. This analysis shows that the only eigenfunctions of $\mc{L}$ among finite order Laurent series in $t$ are those that are members of $\ol S \ker \Gamma$ with eigenvalue $\lambda = 1$.
\begin{example} Let $\Gamma = LE$ in the one-dimensional problem with small steps on the non-negative line with nontrivial operator $E$. Then $\ker \Gamma = \ker E$ consists of functions $\alpha(t) \ol x$, and the eigenfunctions of $\mc{L}$ are $\alpha(t)\ol x \ol S$ for various $\alpha(t)$. In our case, we only look at finite order Laurent series of $\alpha(t)$. Note that (?)
\begin{align}
	Q = \frac{t(\ol x - \ol \sigma)}{(1-tS)(b-a\ol \sigma t)} = \frac{1}{xS}\frac{\sigma - x}{b\sigma - at}\frac{tS}{1-tS}
\end{align}
where we may write
\begin{align}
	\frac{\sigma - x}{b\sigma - at} = \alpha(t) + \beta(t)x
\end{align}
where $\alpha(t)$ and $\beta(t)$ are algebraic functions of $t$ and $a$ since $\sigma = \sigma(t)$. This at least suggests the possibility of connecting the eigenfunctions of $\mc L$ to the solution $Q$.
\end{example}

If we are to have a theory similar to Sturm-Liouville theory, we still need to find a suitable inner product on multivariate Laurent series that will allow us to write solutions for $Q$ in terms of those eigenfunctions with eigenvalue 1.

\section{BLah}
We can go further. We have the following theorem.

\begin{theorem} The full generating function $Q$ for a $d$-dimensional, finite, non-interacting lattice walk model, satisfies a finite-order inhomogeneous PDE in the variables $\row xd$.
\end{theorem}
\begin{proof} We will use the kernel form for the functional equation
\begin{align}
	KQ + t(1- \Gamma)SQ = q.
\end{align}
Here $\Gamma' = 1- \Gamma$ is purely a coefficient extractor and it is finite, since $L$ and $E$ are finite. We will prove the theorem by finding the differential equation explicitly. Since the model is finite, we can bound the degree of $x_k$ that we see on either side of the functional equation. In one scheme we let $A_k$ be the largest power of $x_k$ that we need to clear a negative power of $x_k$ in $\Gamma'$ and in $q$. Then we let $B_k$ be the largest power of $x_k$ we see on either side. Once we multiply by $x_k^{A_k}$, we see that the degree of $x_k$ in the equation is $A_k+B_k$. We take the functional equation and multiply both sides by $x_{[d]}^{A}$. That is,
\begin{align}
	x_{[d]}^{A}(KQ + t\Gamma'SQ) = x_{[d]}^{A}q
\end{align}
and we attain a polynomial equation in the variables $\row xd$. We take the derivative operator
\begin{align}
	D = \del_{[d]}^{A+B+\one} = \del_1^{A_1+B_1+1}\cdots\del_d^{A_d+B_d+1}
\end{align}
and apply it to both sides of the equation to retrieve
\begin{align}
	D(x_{[d]}^{A}(KQ + t\Gamma'SQ)  = D(x_{[d]}KQ) = D(x_{[d]}^{A}q) = 0\,,
\end{align}
which follows since we are applying a larger number of derivatives than the degrees of the monomials and this clears all of them. The other scheme lets us obtain a potentially lower degree differential equation at the cost of homogeneity, if the degrees of $q$ are large. Let $A_k$ be the power of $x_k$ that we need to clear denominators on the left-hand side of the functional equation and $B_k$ be the largest power of $x_k$ on that side. Then we have $D$ as before, but this time it may have lower $A$ and $B$ values, since we are not taking $q$ into account. We have
\begin{align}
	D(x_{[d]}^{A}(KQ + t\Gamma'SQ)  = D(x_{[d]}KQ) = D(x_{[d]}^{A}q)\,.
\end{align}
\end{proof}
\begin{example} For example, let the powers of $x_k$ be $\{-3,-1,4\}$ and $\{-2,0,6\}$ on the left-hand side and right-hand side respectively. We see that $A_k = 3$ and $B_k = 6$. Then we are taking the $A_k + B_k + 1 = 10$-th derivative of both sides of the functional equation. In the second scheme we take $A_k = 3$ still, but $B_k = 4$ and we are taking the 8-th derivative with respect to $x_k$ this time. 
\end{example}

Even in the quarter-plane, we have models with non D-finite generating function (in the $t$ variable), but it turns out that every model so far considered is D-finite in the variables $\row xd$. In order to solve the differential equation we need to investigate the kernel of the differential operator.
\begin{corollary}
The generating function $Q$ of a finite lattice walk model is given by
\begin{align}
	x_{[d]}^{A}KQ(\row xd;t) = \sum_{i=1}^d \sum_{j=0}^{B_i}Aij(x_1,\ldots,x_{i-1},x_i,\ldots, x_d,t) x_i^j\,,
\end{align}
where $A_{ij}$ are functions that need to be determined.
\end{corollary}
This kind of solution would have also appeared if we went through the calculations in the functional equation and extracted coefficients. However, this method is more efficient. If we had used the second scheme in the proof, we would have ended up with a slightly different form but maybe it would have had a slightly lower degree. It would involve an indefinite integral of $q$ and elements in the kernel of the differential operator.

%=======================================
\subsection{Walks With Unconstrained Steps}
%=======================================
We may go through the process again, but this time we have
\begin{align}
	Q = \frac{W}{x_{[d]}^N K}
\end{align}
where $N_k$ is the maximum negative size of steps in the set $\mc{S}$ in the direction $x_k$ or in other words, the maximum $N_k$ for which $S_k^{N_k}$ is nonzero. Then the operators that keep the walks in the quarter plane are
\begin{align}
	E_k = 1 - \sum_{i=1}^{N_k} \ext{\ol x_k^i}\,.
\end{align}
Their conjugates are
\begin{align}
	E_k' = x_k^{N_k} E_k \ol x_k^{N_k} = 1 - \sum_{i=0}^{N_k-1}\ext{x_k^i} = 1 -\mc{E}_k\,,
\end{align}
so in the kernel equation we have
\begin{align}
	x_{[d]}^N = x_{[d]}^N E \ol x_{[d]}^N W = E' W = W +  \sum_{\emptyset \neq I\subseteq [d]} (-1)^{|I|} \mc{E}_I(W)\,.
\end{align}
Here we have various powers $x_k^i$ along with terms with $x_k = 0$ for $k \in I$. Hence, if we seek an operator $D$ that does the same job as the orbit summation operator, we need $D$ to annihilate all the terms with $x_k$'s missing in $W$ on the right-hand side. So then we have
\begin{align}
	D(x_{[d]}KQ) = D(W) = D\left(x_{[d]}^N-\sum_{\emptyset \neq I\subseteq [d]} (-1)^{|I|} \mc{E}_k(W)\right) = D(x_{[d]}^N)\,.
\end{align}
In other words we want $D$ to behave (in some sense) similarly to the derivative
\begin{align}
\del_{[d]}^N = \del_1^{N_1}\cdots \del_d^{N_d}\,.
\end{align}
This suggests some definitions.
\begin{definition}For a sequence $N = (\row Nd)$ of positive integers, we say that a model $\mc{M}$ is of type $N$ if the longest step in $\mc{S}$ in the negative $k$-th direction has length $N_k$ in that coordinate. In particular, we may say that a model is of type $n$ if $N_k = n$ for all $k$.
\end{definition}
For example, the model for full lattice walks with small steps is of type 1.
\begin{definition}For sequences $N = (\row Nd)$ and $M = (\row Md)$ of positive integers, we say that a model $\mc{M}$ is of type $(N,M)$ if it is of type $N$ and the longest step in $\mc{S}$ in the positive $k$-th direction has length $N_k$ in that coordinate. In particular, we may say that a model is of type $(n,m)$ if it is of type $n$, and $M_k = m$ for all $k$.
\end{definition}
For example, full lattice walks with small steps are of type $1,1$. The notation is still a little awful. Anyway, it allows us to write things in the following form
\begin{theorem} A lattice walk model of type $N$ satisfies the functional equation
\begin{align}
	E'W = \fullx xdN
\end{align}
where $E' = \fullx xdN E \fullox xdN$ and $W  = \fullx xdN KQ$.
\end{theorem}

\begin{corollary} For a lattice walk model of type $N$, the generating function $Q$ for the non-interacting case satisfies the PDE
\begin{align}
	\fullx{\del}dN\left(\fullx xdN KQ\right) = 1
\end{align}
and is therefore, D-finite in the $x$ variables.
\end{corollary}
Here $\fullx{\del}dN$ is an example of an operator $D$ that eliminates all the unnecessary terms but does not eliminate $\mon xd$. It is possible to discuss weighted operators in this setting too, but the situation becomes more complicated. However, we may investigate type $(N,M)$ problems by finding a suitable model whose orbital operator is nice enough. This will show that all type $(N,M)$ satisfy an orbital equation with that particular operator, not just a PDE. 

\begin{example} Consider the lattice walk model of type $(N,N)$ with $S = \sum_{k=1}^d \ol x_k^{n_k} + x_k^{n_k}$. Its group $G$ is generated by the elements of groups $G_k$ with $2n_k$ elements each - each of the elements $g_{k,j,l} = \omega_{n_k}^jx_k^l$ where $\omega_{n_k} = \exp(2\pi i/n_k)$ and $l\in\{-1,1\}$. These elements commute and allow us to find an orbital operator that annihilates all terms on the left-hand side except  $\fullx xdN KQ$ and also spares the $\fullx xdN$ term on the right-hand side.
\end{example}



%=======================================
\section{Interaction with $\del_t$}
%=======================================
Here we'll investigate $\del_t$ interacts with $\mc{L}$.
\begin{align}
	\del_t \mc{L} &= \del_t - \del_t(t\Gamma S) = \del_t - \Gamma S- t\Gamma S\del_t = \mc{L}\del_t - \Gamma S\,,\\
	\del_t^2 \mc{L} &= \mc{L}\del_t^2 - 2\del_t \Gamma S\,,\\
	\del_t^n \mc{L} &=\mc{L} \del_t^n - n\del_t^{n-1}\Gamma S
\end{align}
for integer $n$. Another way to write these is with bracket notation
\begin{align}
	[\mc{L},\del_t^n] = \mc{L}\del_t^n-\del_t^n\mc{L} = n\del_t^{n-1}\Gamma S\,.
\end{align}

%=======================================
\section{Interaction with $\del_i$}
%=======================================
Let's take the operator $\Gamma_n(a) = 1 + b x^n[x^n]$ again. Let $\del_k = \del/\del_k$ and in this case just $\del$. We have
\begin{align}
	\del \Gamma_n(a) = \del + b\del(x^n[x^n]) = \del + bnx^{n-1}[x^n]\,.
\end{align}
Now if $n\ge 0$ we have
\begin{align}
	\del^{n+1}\Gamma_n(a)  = \del^{n+1} + b\del^{n+1}(x^n)[x^n] = \del^{n+1}\,.
\end{align}
If $n\leq 0$ taking derivatives will yield ever greater powers of $\ol{x}$ so we want to pre-multiply by a power of $x$ first.
\begin{align}
	\del x^{-n} \Gamma_{n}(a) = \del x^{-n} + b\del [x^n] = \del x^{-n} = -nx^{-n-1} + x^{-n}\del\,.
\end{align}
If we have more terms in an operator then we need to multiply by a power of $x$ to clear out any negative powers of $x$. Then we take the lowest order derivative in order to clear out all the extraction terms. Suppose that
\begin{align}
	\Gamma = 1 + \sum_{i=0}^{m} a_i \ext{x^i} - \sum_{i=1}^n b_i\ext{\ol{x}^i}
\end{align}
where $m$ and $n$ are non-negative. We need to pre-multiply by $x^n$ to clear out all the negative powers of $x$. Then we have
\begin{align}
	x^n\Gamma = x^n + \sum_{i=0}^{m} a_i \ext{x^{i+n}} - \sum_{i=1}^n b_i\ext{x^{n-i}}\,.
\end{align}
This operator is cleared by $\del^{m+n+1}$.


\endinput
